\scanpage{128}

% manuscript lines / position: 1–2
\[
 \begin{aligned}
 \tilde b_i&=\varpi_i-\chi q_i,\qquad b_i+\tilde b_i=\chi\mathrm{id}-\chi\mathrm{id}=0\\
 &=\underline{\underline{-b_i}}
 \end{aligned}
\]

% manuscript lines / position: 3–4
\[
 \begin{aligned}
 \tilde U_i\ \text{tel que}\quad q'_i&=\underbrace{q_i}_{\mathrm{id}-p_i}-\tilde U_i
 =\deleted{\unsure{\chi}}\mathrm{id}-\underbrace{(p_i+\tilde U_i)}_{p'_i}=\unsure{\cdot}\\
 \tilde U_i&=-U_i
 \end{aligned}
\]

% manuscript lines / position: 5–6
\[
 \begin{aligned}
 \deleted{\unsure{U}}\tilde U&=\tilde b_1\otimes\tilde U_2+\tilde U_1\otimes\tilde b_2+\chi\tilde U_1\otimes\tilde U_2\\
 &=b_1\otimes U_2+U_1\otimes b_2+\chi U_1\otimes U_2=U
 \end{aligned}
\]

% manuscript lines / position: 5–6, right of vertical separator
\margin{Mais Attention, $\tilde T_{\tilde U}=\deleted{\unsure{+}}T_{-\tilde U}=-T_{-U}=T_U^{-1}$}

% manuscript lines / position: 7
\[
 \tilde b=\tilde b_1\tilde b_2=b_1b_2
\]

% manuscript lines / position: 8–11, left
\diagram{\scriptsize\begin{tikzcd}[ampersand replacement=\&,row sep=large]
 {\unsure{\mathcal E|\mathscr E}_{(p_1,q_1)(p_2,q_2)}\deleted{\unsure{\simeq}}=M_1\otimes M_2\unsure{+|\oplus}L_1\otimes L_2}
 \arrow[d,"{T_{v_{(p_1,q_1)(p_2,q_2)}}}","{\sim}"']\\
 {\mathcal E_{(\tilde p_1,\tilde q_1)(\tilde p_2,\tilde q_2)}=M_1\otimes M_2\unsure{+|\oplus}L_1\otimes L_2}
\end{tikzcd}}

% manuscript lines / position: 8–9, right of upper object
$\chi$-scindage défini \inserted{(conv. \unsure{g.|g})} par

% manuscript lines / position: 9, upper right
\[
 M_1\otimes M_2\xrightarrow{b=b_1b_2}L_1\otimes L_2
\]

% manuscript lines / position: 10–11, right of lower object
$\chi$ sc. défini (conv. \unsure{g.}) par $-b$

% manuscript lines / position: 12–14
$T_v$ \inserted{$(=(y,z)\longmapsto y,z\deleted{\unsure{-}}\unsure{+}v(y))$} \deleted{\unsure{tr}} transforme le $\chi$-scindage défini par $b$ en le $\unsure{\phi|\chi}$-scindage défini \deleted{par} $b+\chi v$.

% manuscript lines / position: 15
\[
 \deleted{\unsure{\text{Si}}}\ \text{On veut que}\quad b+\chi v=\underset{\substack{\Vert\\-b}}{\tilde b}\deleted{\unsure{=}}
 \quad\text{i.e.}\quad v\overset{!}{=}-\frac2{\chi}b
\]

% manuscript lines / position: 16–17
Supposons que $\exists\,\gamma\in k$, tq. $\gamma\chi=2$, on peut prendre \deleted{\unsure{$\xi$}}

% manuscript lines / position: 17
\[
 \boxed{v=-\gamma b}
\]

% manuscript lines / position: bottom left
% archive sheet number: 127
